\documentclass[12pt,a4paper]{article}

\usepackage[T2A]{fontenc}
\usepackage[utf8]{inputenc}
\usepackage[russian]{babel}
\usepackage[a4paper,left=30mm,right=15mm,top=20mm,bottom=20mm]{geometry}
\usepackage{amsmath}
\usepackage{array}
\usepackage{booktabs}
\usepackage{listings}
\usepackage[hidelinks]{hyperref}

\lstset{
  basicstyle=\ttfamily\small,
  breaklines=true,
  columns=fullflexible,
  keepspaces=true,
  frame=single,
  tabsize=4,
  showstringspaces=false
}

\begin{document}

\begin{center}
  \textbf{КОНТРОЛЬНАЯ РАБОТА}\\[0.5em]
  по дисциплине «Базы данных»\\[1em]
  \textbf{Вариант 19}
\end{center}

\vspace{1em}
\noindent Выполнил: Лузан Роман\\[0.8em]
\noindent Группа: в5130904/40321

\section*{Вариант 19}
Определить номера проектов, в которых используется по крайней мере одна деталь, имеющаяся у поставщика с номером \texttt{S1}.

\section*{Обозначения}
Используется классическая схема поставок:
\begin{itemize}
  \item $S(S\#, SNAME, STATUS, CITY)$ --- поставщики;
  \item $P(P\#, PNAME, COLOR, WEIGHT, CITY)$ --- детали;
  \item $J(J\#, JNAME, CITY)$ --- проекты;
  \item $SPJ(S\#, P\#, J\#, QTY)$ --- поставки деталей поставщиками для проектов.
\end{itemize}
Атрибуты $S\#$, $P\#$ и $J\#$ являются номерами поставщика, детали и проекта соответственно. Имена проектов хранятся в отношении $J$ в атрибуте $JNAME$.

\section*{Реляционная алгебра}
Обозначим две копии отношения поставок через $SPJ_1$ и $SPJ_2$. Первая копия задаёт детали, имеющиеся у поставщика \texttt{S1}; вторая связывает эти детали с проектами, где они используются. Проекция возвращает номера проектов без повторов:
\[
\pi_{SPJ_2.J\#}\Bigl(\sigma_{SPJ_1.S\#='S1' \land SPJ_1.P\#=SPJ_2.P\#}(\rho_{SPJ_1}(SPJ) \times \rho_{SPJ_2}(SPJ))\Bigr).
\]

\section*{Реляционное исчисление}
Номер проекта включается в результат, если существует поставка детали этому проекту и существует поставка той же детали от поставщика \texttt{S1}:
\[
\{\, t.J\# \mid SPJ(t) \land \exists u\,(SPJ(u) \land u.S\#='S1' \land u.P\#=t.P\#)\,\}.
\]

\section*{Запрос SQL}
\begin{lstlisting}[language=SQL]
SELECT DISTINCT project_supply."J#"
FROM "SPJ" AS project_supply
WHERE EXISTS (
    SELECT 1
    FROM "SPJ" AS s1
    WHERE s1."S#" = 'S1'
      AND s1."P#" = project_supply."P#"
)
ORDER BY project_supply."J#";
\end{lstlisting}
Условие \texttt{EXISTS} проверяет, что деталь из поставки проекта также присутствует среди деталей, поставляемых \texttt{S1}. \texttt{DISTINCT} оставляет каждый номер проекта в результате один раз.

\section*{Результат}
В стандартном учебном наборе данных детали поставщика \texttt{S1} используются в проектах:
\begin{center}
\begin{tabular}{c}
\toprule
Номер проекта \\
\midrule
J1 \\
J4 \\
\bottomrule
\end{tabular}
\end{center}

\section*{Ответ}
Номера проектов, в которых используется хотя бы одна деталь, имеющаяся у поставщика \texttt{S1}: \textbf{J1, J4}.

\end{document}